\begin{frame}{Uma função transforma entradas em resultado}
\begin{center}
\begin{tikzpicture}[node distance=0.5cm, every node/.style={align=center,font=\small}]
\node[draw,rounded corners,text width=2.2cm] (entrada) {Temperatura: 32\\Limite: 30};
\node[draw,fill=green!10,right=of entrada,text width=3cm] (processo) {\texttt{alarme\_ativo}\\comparar com o limite};
\node[draw,rounded corners,right=of processo,text width=2.1cm] (saida) {Resultado\\\texttt{true}};
\draw[-{Latex}] (entrada)--(processo);
\draw[-{Latex}] (processo)--(saida);
\end{tikzpicture}
\end{center}
\begin{block}{Separar decisão de ação}
A função calcula se há alarme. Outra parte do programa imprime a mensagem ou, futuramente, aciona uma saída digital.
\end{block}
\small Assim, podemos testar a regra sem conectar o equipamento.
\end{frame}
